\section{Related Works}
\label{sec:formatting}

%-------------------------------------------------------------------------
\subsection{Audio-Driven 3D Talking Head Synthesis}
With the emergence of NeRF~\cite{mildenhall2021nerf} and 3D Gaussian Splatting (3DGS)~\cite{kerbl20233d}, 3D talking head synthesis has become a promising direction. Existing 3D person-specific methods ~\cite{li2024talkinggaussian, zhang2025fate, li2025rgbavatar, hou2022face, hou2021towards} train a model for each identity from scratch, achieving high-fidelity results but requiring several minutes of video. To reduce this cost, one-shot methods~\cite{ye2024real3d, hu2025ggtalker, tan2025fixtalk} adapt pretrained priors from a single portrait, but their limited identity information often yields less personalized and unstable results. Alternatively, few-shot Pretrain-and-Adapt methods such as InsTaG~\cite{li2025instag} and FIAG~\cite{nie2025few} adapt a universal motion prior learned from a multi-identity corpus to new identities using only a few seconds of video, enabling efficient and realistic personalization. However, these few-shot methods face two key limitations: the lack of explicit emotional modeling, leading to less expressive facial animations; and unconstrained 3DGS deformation, which causes geometric instability under diverse and exaggerated expressions. Our work addresses both limitations by introducing emotion-aware motion learning with a strong structural prior.

%-------------------------------------------------------------------------
\subsection{Emotional Talking Head Synthesis}
Natural talking head synthesis requires not only accurate lip synchronization but also emotionally expressive facial dynamics. Recent 2D generative approaches~\cite{tan2023emmn, tan2024edtalk, zakharov2019few, xu2024hallo, tan2024flowvqtalker, tan2025disentangle, tan2025edtalk++} produce vivid expressions, but the absence of explicit 3D structural priors makes it difficult to preserve consistent geometry and realistic deformation, particularly under large head poses or highly expressive facial motions. In the 3D domain, EMOTE~\cite{danvevcek2023emotional} and EmoVOCA~\cite{nocentini2025emovoca} rely on 3D face datasets with manual emotion annotations to train FLAME-based models. However, such annotation-driven pipelines are restricted by coarse discrete emotion categories and annotator subjectivity, preventing them from learning fine-grained emotional dynamics directly from audio. EmoTalk3D~\cite{he2024emotalk3d} achieves high-fidelity results with 3D Gaussian Splatting, but its reliance on person-specific optimization significantly limits scalability for rapid personalization. To our knowledge, no methods simultaneously achieve few-shot adaptation, 3D consistency, and emotion-aware modeling from audio, as achieved by EmoTaG.
